\subsection{Python Compilers, Accelerators, and Translators: Cython, Nuitka, Numba, MyPyC, and Related Tools}

After separating Python from CPython and after comparing several Python implementations, another category must be separated: compilers, accelerators, translators, and packaging compilers. These tools are often discussed as if they all answer the same question: “How do I make Python faster?” In reality, they answer different questions.

Some tools translate Python-like code into C extension modules. Some compile selected functions at runtime. Some turn Python programs into native executables that still carry a Python runtime. Some use type annotations to generate more efficient C-level code. Some are mainly distribution tools, not performance tools. If these categories are mixed together, the student receives a false picture of Python compilation.

The central fact is this: ordinary Python is hard to compile like C because ordinary Python deliberately leaves many decisions until runtime. A name can be rebound to a different kind of object. An operator can call user-defined methods. Attribute access can trigger descriptors. Imports can execute code. Classes can be modified. Functions are objects. The program can inspect and sometimes modify its own execution environment. This flexibility is useful, but it prevents a simple ahead-of-time compiler from treating Python source as if it were statically typed C source.

Therefore, most successful Python acceleration tools do not compile all possible Python behavior equally. They restrict part of the language, specialize part of the program, use type information, delegate work to native libraries, or keep CPython nearby.

\subsubsection{Why Python Compilation Is Not C Compilation}

In C, many facts are known before the program runs. The compiler knows that an \texttt{int} variable stores an integer of a certain representation. It knows the layout of structures. It knows function signatures. It can generate direct machine instructions for arithmetic, memory access, and function calls.

In ordinary Python, this is not generally true. Consider:

\begin{verbatim}
x = a + b
\end{verbatim}

At compile time, CPython may know that the source contains a name load, another name load, an addition operation, and a store. But it does not necessarily know what objects \texttt{a} and \texttt{b} will refer to when the line executes. They may be integers, floats, strings, lists, NumPy arrays, user-defined objects, or objects whose \texttt{\_\_add\_\_()} method has unusual behavior.

So the compiler cannot simply emit one CPU addition instruction in the general case. It must preserve Python's runtime semantics. This means object lookup, type dispatch, method resolution, exception handling, reference management, and fallback behavior.

This is why Python acceleration usually depends on one of four strategies:

\begin{itemize}
    \item add type information;
    \item restrict the accepted Python subset;
    \item compile only selected hot functions;
    \item keep the Python runtime but reduce overhead around it.
\end{itemize}

\subsubsection{Cython: Python Near the C Boundary}

Cython is best understood as a bridge between Python and C. It can compile Python-like code into C or C++ code that becomes a CPython extension module. It also allows the programmer to add C-level type declarations, call C functions directly, and define data structures closer to native representation.

A simple Cython file may look almost like Python. But the performance gain becomes stronger when the programmer gives Cython more information:

\begin{verbatim}
def add_python(a, b):
    return a + b

cdef int add_c(int a, int b):
    return a + b
\end{verbatim}

The first function still behaves like ordinary Python. The second function gives Cython C-level information. Now the operation can be compiled much closer to ordinary C integer addition.

Cython is therefore not magic. It becomes powerful when the programmer deliberately moves selected parts of the program toward C. This is especially useful for tight loops, numeric kernels, wrappers around existing C libraries, and extension modules that must expose a Python interface while doing heavy work in native code.

Its cost is that the code is no longer purely ordinary Python in spirit. The programmer begins to think about C types, C arrays, memoryviews, extension-module building, compilation, and native binary compatibility. Cython gives speed by partially abandoning the fully dynamic surface of Python where necessary.

\subsubsection{Nuitka: Whole-Program Compilation with Runtime Preservation}

Nuitka has a different goal. It compiles Python programs into C/C++ code and then into executables or extension modules. But this should not be misunderstood as “Python becomes C in the same sense that C source becomes C machine code.” Nuitka must still preserve Python semantics.

This means that a compiled Nuitka program may still carry or depend on the Python runtime machinery needed to execute Python behavior correctly. The output can be useful for distribution, startup control, packaging, and sometimes performance, but the dynamic nature of Python does not disappear simply because the program was compiled.

Nuitka is useful when one wants to ship a Python application as a compiled artifact, reduce some interpreter overhead, or produce a more self-contained program. It is not the same kind of tool as Cython. Cython asks the programmer to expose C-level facts in selected code. Nuitka tries to accept normal Python programs and compile them while preserving Python behavior.

The distinction is important:

\begin{itemize}
    \item Cython is often used for extension modules and native-speed inner loops;
    \item Nuitka is often used for compiling larger Python programs or packaging them as executables.
\end{itemize}

Both involve compilation, but they sit at different architectural points.

\subsubsection{Numba: Runtime Compilation of Numerical Kernels}

Numba is mostly useful for numerical functions. Its common pattern is simple: the programmer writes a Python function, decorates it, and Numba compiles suitable versions of that function at runtime.

A typical example is:

\begin{verbatim}
from numba import njit

@njit
def sum_values(values):
    total = 0.0
    for x in values:
        total += x
    return total
\end{verbatim}

The important detail is that Numba does not try to compile every possible Python program equally. It works best when the function uses numeric types, arrays, loops, and operations that can be lowered into efficient machine code. In such cases, Numba can avoid much of the ordinary CPython bytecode dispatch overhead.

Numba's strength is that it allows a programmer to keep a Python-like development style while accelerating selected hot functions. Its weakness is that it is not a universal Python compiler. Code that depends heavily on arbitrary Python objects, dynamic attribute tricks, complex class behavior, or unsupported library calls may fall outside the efficient compilation path.

Architecturally, Numba is very different from Cython. Cython is normally an ahead-of-time translation path through C/C++ extension modules. Numba is normally a just-in-time specialization path for selected functions, especially numerical functions.

\subsubsection{MyPyC: Type-Annotated Python Compiled to C Extensions}

MyPyC uses Python type annotations to compile modules into CPython C extension modules. Its central idea is that type hints can be used not only for checking, but also for producing faster code.

For example, ordinary Python type hints such as:

\begin{verbatim}
def add(a: int, b: int) -> int:
    return a + b
\end{verbatim}

give the compiler more information than an unannotated function. If the tool can rely on those types under its rules, it can generate code with less dynamic overhead.

MyPyC is especially interesting because it grows from modern Python's gradual typing culture. Type hints were not originally added as a C-like compilation system. They were mainly introduced for readability, tooling, checking, and large-codebase maintainability. But once type information exists, compilation becomes more plausible for selected code.

The tradeoff is that MyPyC does not preserve the full freedom of arbitrary Python in every corner. It works best with code that is already disciplined, typed, and written in a style compatible with its restrictions. This makes it attractive for accelerating large typed Python codebases, but not for every dynamic script.

\subsubsection{Packaging Tools Are Not Necessarily Compilers}

Some tools are often confused with compilers because they produce executable-looking artifacts. Examples include tools that bundle a Python program together with an interpreter, libraries, and resources. Such tools may create an \texttt{.exe} or an application bundle, but that does not necessarily mean the Python code has been translated into optimized native machine code.

This distinction matters. A packaged Python program may still execute Python bytecode through an embedded interpreter. The user sees a native-looking file, but internally the architecture may still be:

\begin{verbatim}
application executable
    contains or launches Python runtime
        loads Python bytecode or source
            executes through interpreter
\end{verbatim}

This is useful for deployment. It is not automatically a performance transformation.

\subsubsection{The Common Pattern: Move the Hot Path Downward}

The practical reason these tools exist is that most programs are not uniformly expensive. A large application may spend most of its time in a few small regions: a numerical loop, a parser, an image operation, a geometry transformation, a serialization step, or a repeated calculation.

Python acceleration often means identifying those hot paths and moving them downward:

\begin{itemize}
    \item from Python objects to typed native values;
    \item from interpreter loops to compiled loops;
    \item from Python lists to contiguous arrays;
    \item from dynamic dispatch to known operations;
    \item from pure Python functions to native extension modules.
\end{itemize}

This preserves Python's role as the orchestration layer while moving the expensive inner work closer to the machine.

This is exactly the same architectural reason Python became dominant in scientific computing. Python remains the outer language of expression, experiment, and coordination. The inner loops move to C, C++, Fortran, LLVM-generated code, GPU kernels, or specialized native libraries.

\subsubsection{A Practical Comparison}

The following simplified comparison is useful:

\begin{itemize}
    \item \textbf{Cython} is best when the programmer is willing to add C-level information and build extension modules.
    \item \textbf{Nuitka} is best when the goal is compiling or packaging larger Python programs while preserving Python semantics.
    \item \textbf{Numba} is best for accelerating selected numerical functions at runtime.
    \item \textbf{MyPyC} is best for compiling disciplined, type-annotated Python modules into CPython extensions.
    \item \textbf{Bundling tools} are best for distribution, not necessarily for speed.
\end{itemize}

These tools do not replace one another. They solve different problems.

\subsubsection{The Conceptual Lesson}

The existence of these tools proves that Python is not a single execution story. Python code can be interpreted, JIT-compiled, translated to C, compiled into extension modules, bundled with an interpreter, or used as a control layer above native libraries.

But the same warning must be repeated: a tool that accelerates Python does not necessarily turn Python into C. Python's full language semantics are too dynamic for that simple statement. Speed usually appears when the tool can reduce uncertainty. Type declarations, type hints, numerical arrays, restricted subsets, and stable hot loops all give the compiler more facts. More facts allow less runtime dispatch. Less runtime dispatch allows faster execution.

For this course, the architectural conclusion is direct. Python's surface simplicity is supported by heavy runtime machinery. Acceleration tools work by bypassing, specializing, or reducing parts of that machinery in selected regions. They do not abolish the CPython model; they exploit boundaries inside it.

Thus, Python performance should not be understood as one number. There is pure Python interpreter speed. There is CPython extension speed. There is vectorized native-library speed. There is JIT-specialized speed. There is packaged-application behavior. Serious Python programming requires knowing which layer is currently doing the work.